\chapter{Verilog 연산자의 실제 자원 비용}

\section{연산자 하나를 같은 조건에서 합성해 본다}

Verilog의 짧은 기호와 FPGA primitive 수는 일대일로 대응하지 않는다. 이를
확인하기 위해 23-bit datapath에서 자주 쓰는 연산을
\rtlfile{operator\_cost\_all} top 하나에 각각 독립 instance로 배치하고 한 번의
OOC synthesis를 수행했다.

\begin{itemize}
  \item Vivado 2020.2, \rtlfile{xc7z020clg484-1}
  \item \rtlfile{flatten\_hierarchy none}
  \item arithmetic instance의 input/output 사이에 한 register boundary
  \item instance 사이의 흡수와 제거를 막기 위한 hierarchy 보존
  \item place/route 이전의 post-synthesis primitive count
\end{itemize}

\begin{longtable}{L{41mm}r r r r L{46mm}}
\toprule
RTL 연산 & LUT & CARRY4 & FF & DSP & 해석 \\
\midrule
23-bit \rtlfile{a+b} & 23 & 6 & 69 & 0 & 한 carry chain \\
23-bit \rtlfile{a-b} & 23 & 6 & 69 & 0 & 한 carry chain \\
23-bit \rtlfile{a==b} & 8 & 2 & 47 & 0 & bit 비교와 reduction \\
23-bit unsigned \rtlfile{a<b} & 12 & 3 & 47 & 0 & LUT packing된 magnitude compare \\
\rtlfile{a>=8380417} & 6 & 0 & 24 & 0 & constant bit pattern으로 축소 \\
23-bit 2:1 MUX & 23 & 0 & 70 & 0 & output bit당 LUT 하나 \\
23-bit 4:1 MUX & 23 & 0 & 117 & 0 & 4 data+2 select가 LUT6에 들어감 \\
23-bit constant left shift & 0 & 0 & 41 & 0 & 배선과 constant bit \\
23-bit variable left shift & 64 & 0 & 51 & 0 & multi-level barrel MUX \\
23$\times$16 multiply & 0 & 0 & 39 & 1 & DSP48E1로 흡수 \\
23-bit modular subtraction & 46 & 12 & 92 & 0 & subtract와 conditional add \\
13-input reduction OR & 3 & 0 & 14 & 0 & 작은 reduction tree \\
\bottomrule
\end{longtable}

\newpage
같은 top에 storage inference 예제도 넣었다.

\begin{longtable}{L{48mm}r r r r r L{37mm}}
\toprule
RTL 구조 & CLB LUT & LUTRAM & SRL & FF & R36 & 해석 \\
\midrule
16-cycle $\times$ 23-bit delay, no reset & 23 & 0 & 23 & 46 & 0 & SRL16E와 입출력 FF \\
16-cycle $\times$ 23-bit delay, reset & 0 & 0 & 0 & 368 & 0 & $16\times23$ FF chain \\
32$\times$8 asynchronous-read RAM & 8 & 8 & 0 & 0 & 0 & distributed LUTRAM \\
576$\times$23 synchronous RAM & 0 & 0 & 0 & 0 & 1 & RAMB36E1 \\
\bottomrule
\end{longtable}

FF 수에는 연산기 자체뿐 아니라 실험 조건을 고정한 input/output boundary
register가 포함된다. 예를 들어 23-bit add의 69 FF는 두 23-bit input register와
23-bit output register다. 그러므로 표의 FF 열을 순수 operator cost로 읽지
않는다. 반면 LUT, CARRY4와 DSP 열은 각 독립 instance의 조합 구조를 직접
보여 준다.

\section{Addition과 subtraction: LUT가 carry를 구동한다}

7-series의 CARRY4 하나는 네 bit의 carry propagation을 담당한다. 이 실험에서
23-bit addition과 subtraction은 각각
\[
\left\lceil\frac{23}{4}\right\rceil=6
\]
개의 CARRY4와 23 LUT를 사용했다. LUT는 각 bit의 propagate/generate 또는
sum 선택을 만들고, CARRY4가 인접 bit의 carry를 빠르게 전달한다.

이 결과가 “23-bit adder는 항상 23 LUT”라는 법칙은 아니다. constant operand,
사용되지 않는 상위 carry, 주변 MUX와 downstream logic이 있으면 합성기가
LUT를 합치거나 제거한다. 그러나 RTL을 처음 검토할 때
\textbf{23-bit carry 연산 하나는 대략 여섯 CARRY4 길이}라고 보는 것은 유용하다.

\subsection{왜 modular subtraction이 무거운가}

한 cycle modular subtraction의 실험 코드는 raw difference와 conditional
addition을 연속으로 둔다.

\begin{lstlisting}
diff      = {1'b0,a} - {1'b0,b};
corrected = {1'b0,diff[22:0]} + {1'b0,q};
result    = diff[23] ? corrected[22:0] : diff[22:0];
\end{lstlisting}

합성 결과는 46 LUT와 12 CARRY4였다. 일반 subtract의 정확히 두 배인 carry
resource가 한 register 경계에 놓인 셈이다. 150~MHz 실패 path에서 관찰한
11 CARRY4도 이 구조와 같은 직관을 준다. 최종 RTL이 early difference를
raw subtract와 conditional correction의 두 cycle로 나눈 이유를 area와
timing 양쪽에서 설명할 수 있다.

\section{Comparator: 같은 23 bit라도 비용이 다르다}

\rtlfile{==}는 bit별 equality와 reduction으로 구현할 수 있다. 이 실험의
23-bit equality는 CLB LUT 8개와 CARRY4 2개였다. unsigned magnitude compare는
bit의 우선순위를 판단해야 하며 CLB LUT 12개와 CARRY4 3개를 사용했다. netlist의
개별 LUT function 수와 hierarchy report의 physical CLB LUT 수는 LUT6\_2 packing
때문에 다를 수 있으므로, 이 표는 hierarchy utilization의 CLB LUT를 기준으로 했다.

constant compare는 더 작아질 수 있다. ML-DSA modulus
8,380,417과의 23-bit \rtlfile{>=} 비교는 6 LUT, CARRY4 0개로 줄었다.
constant의 0/1 pattern과 reachable range를 합성기가 이용했기 때문이다.

따라서 다음 두 RTL은 폭이 같아도 같은 면적으로 예상하지 않는다.

\begin{lstlisting}
wire ge_runtime = (value >= runtime_q);
wire ge_mldsa   = (value >= 23'd8380417);
\end{lstlisting}

현재 NTT RTL이 외부의 unconstrained \rtlfile{PRM\_Q}를 넓은 arithmetic cone에
계속 전달하지 않고 scheme별 internal constant를 선택하는 이유도 여기에 있다.
constant propagation은 comparator뿐 아니라 adder의 고정 bit와 MUX input도
줄일 수 있다.

\section{MUX: input 개수와 bit 폭을 함께 센다}

23-bit 2:1 MUX와 4:1 MUX는 이 isolated synthesis에서 모두 23 LUT였다.
한 output bit를 기준으로 2:1 MUX는 두 data와 select의 3-input function이고,
4:1 MUX는 네 data와 두 select의 6-input function이다. 둘 다 LUT6 하나에
들어간다.

이 결과에서 4:1 MUX가 항상 2:1과 같은 timing이라는 결론은 나오지 않는다.
LUT input pin delay, select fanout와 placement가 다르다. 8:1 이상이 되거나
MUX 앞뒤에 arithmetic이 붙으면 LUT level 또는 dedicated MUX resource가
더 필요할 수 있다.

특히 다음 두 표현은 기능이 같아 보여도 합성 문맥이 다르다.

\begin{lstlisting}
// Multiple arithmetic branches may become parallel operators plus a result MUX.
y = mode0 ? add_mod(a,b) :
    mode1 ? add_mod(c,d) : add_mod(e,f);

// Select operands first, then use one physical arithmetic position.
case (mode)
  0: begin op_a=a; op_b=b; end
  1: begin op_a=c; op_b=d; end
  default: begin op_a=e; op_b=f; end
endcase
y = add_mod(op_a,op_b);
\end{lstlisting}

최종 butterfly는 두 번째 구조를 사용한다. mode branch마다 function을 호출해
여러 adder와 result MUX를 만들지 않고, physical position의 operand를 먼저
선택한 뒤 arithmetic function을 한 번 호출한다.

\section{Shift: constant와 variable은 전혀 다른 회로다}

같은 \rtlfile{<<} 연산자라도 shift amount가 compile-time constant인지
runtime signal인지가 중요하다. 23-bit \rtlfile{a<<5}는 LUT를 하나도 쓰지
않았다. bit 위치를 바꾸는 배선과 0 constant로 충분하기 때문이다.

반면 5-bit runtime shift amount를 사용한 23-bit left shift는 64 LUT였다.
여러 shift 후보를 단계별로 고르는 barrel MUX가 필요하다. NTT address와
modular arithmetic에서 고정 shift를 적극적으로 사용하는 이유는 곱셈을
덧셈으로 바꾸는 것뿐 아니라 variable barrel network를 피하기 위해서다.

\section{Multiplication과 reduction}

\rtlfile{23x16} multiply에는 \rtlfile{use\_dsp}를 지정했고 결과는 LUT 0,
DSP48E1 한 개였다. operand가 DSP48E1 multiplier port 범위에 들어가므로
product와 output register가 hard block에 흡수된다. 최종
\rtlfile{unifiedMUL}은 이런 DSP를 두 개 사용하고, 두 instance의 합계가
네 DSP다.

constant multiplication은 항상 DSP를 써야 하는 것은 아니다. 최종
$1006x\bmod3457$ block은 두 6-bit table과 addition/correction으로 구현되어
standalone post-route hierarchy에서 84 LUT와 25 FF, DSP 0개를 사용한다.
세 lane의 modular division by 3 block은 shift와 addition으로 구현되어
465 LUT와 216 FF, DSP 0개다. hard DSP를 네 개로 제한하는 대신 fabric
area를 사용한 명시적인 trade-off다.

13-input reduction OR는 isolated synthesis에서 3 LUT였다. 수만 보면 작지만
그 결과가 \rtlfile{BUSY}, enable과 memory phase MUX를 동시에 구동하면
fanout과 route가 critical path가 될 수 있다. 최종 core가
\rtlfile{|valid\_pipe}를 occupancy counter로 바꾼 이유는 LUT 세 개 자체보다
그 reduction 결과의 timing cone을 없애기 위해서다.

\section{같은 delay가 SRL 또는 368 FF가 되는 이유}

23-bit data를 16 cycle 미루는 두 RTL도 함께 합성했다. reset이 없는 delay는
23개의 SRL16E와 46개의 입출력 FF로 mapping됐다. 같은 delay vector 전체에
synchronous reset을 추가한 구조는 SRL을 사용하지 못하고 정확히
$23\times16=368$개의 FF로 남았다.

이 결과를 보고 내가 얻은 기준은 ``delay가 보이면 SRL''이 아니라 다음과 같다.

\begin{enumerate}
  \item 중간 tap을 독립적으로 읽거나 수정하는가?
  \item data 자체를 reset해야 하는가, valid만 reset하면 되는가?
  \item enable과 delay 길이가 SRL template에 맞는가?
  \item 합성 report에 SRL16E/SRLC32E가 실제로 나타났는가?
\end{enumerate}

reset을 없애기 위해 기능을 바꾸면 안 된다. valid가 0인 동안 data가 관찰되지
않는다는 protocol을 먼저 확인하고 first-valid cycle을 regression으로 검사한다.

\section{같은 저장량도 LUTRAM 또는 BRAM이 된다}

32$\times$8 asynchronous-read memory는 8개의 CLB LUT를 memory로 사용했다.
576$\times$23 synchronous memory는 RAMB36E1 하나가 됐다. 13,248 bit라는
capacity만 보면 RAMB18에 가까워 보이지만 576-word depth와 23-bit width를 함께
지원하는 configuration 때문에 RAMB36 하나로 mapping됐다.

이 경험 때문에 memory area를 계산할 때 단순히 total bit를 BRAM capacity로
나누지 않는다. depth, width, read/write port와 synchronous read latency를 먼저
적고 실제 primitive configuration을 확인한다.

\section{면적을 예상하는 순서}

\begin{enumerate}
  \item 각 cycle 경계 사이의 add/sub carry chain 수와 bit 폭을 센다.
  \item arithmetic 앞 operand MUX와 뒤 result MUX 가운데 어느 쪽이 생기는지
        그린다.
  \item variable shift, runtime modulus, generic multiply처럼 작은 기호가
        큰 network를 만드는 표현을 표시한다.
  \item pipeline data뿐 아니라 valid, mode, address와 destination tag의
        register bit도 센다.
  \item memory는 선언 bit 수가 아니라 target primitive의 width/depth
        configuration으로 올림해 센다.
  \item 마지막에는 반드시 post-synthesis hierarchy와 primitive count로
        예상이 맞는지 확인한다.
\end{enumerate}

이 프로젝트의 최종 standalone hierarchy에서 complete
\rtlfile{unifiedBUT}은 2,704 LUT, 2,038 FF와 네 DSP를 사용한다.
하위 두 \rtlfile{unifiedMUL}은 같은 RTL이지만 각각 459/378 LUT로 다르게
보였다. downstream trimming과 hierarchy 간 LUT combining 때문이다.
하위 block 수를 단순 합산하거나 동일 RTL의 한 instance 수치를 절대값으로
일반화하지 않는다.

\begin{checklistbox}
\begin{itemize}
  \item add/sub는 먼저 $\lceil W/4\rceil$ CARRY4 chain을 예상한다.
  \item conditional correction은 두 번째 carry chain과 MUX가 붙는지 확인한다.
  \item 4:1 bit MUX는 LUT6 하나에 들어가도 select fanout은 별도로 본다.
  \item constant shift는 배선, variable shift는 barrel MUX로 생각한다.
  \item multiplication은 DSP inference report와 내부 register 흡수를 확인한다.
  \item 작은 OOC 결과와 complete-system post-route 결과를 같은 수치로 취급하지 않는다.
\end{itemize}
\end{checklistbox}

재현용 RTL, Tcl과 실제 summary는 repository의
\rtlfile{examples/operator\_cost}에 있다. 하나의 in-memory project에서
synthesis를 한 번만 실행하므로 Vivado project directory나 checkpoint를
저장소에 남기지 않는다.
